\documentclass{subfiles}
\begin{document}
    \begin{procedure}
        \caption{Diffie-Hellman-Key-Exchange()}\label{proc:DH-exchange}
        \nl Let $\T$ be a trusted third party. 
            $\T$ chooses a prime number $p$ and chooses an element $\gamma \in \gf_{p}$.
            \DontPrintSemicolon\;

        \nl Consider $\B$; it chooses an integer $b$ and computes $\beta \equiv \gamma^{b} \mod$,
            which would represent its public key.
            In the same manner $\A$ computes its public key $\alpha \equiv \gamma^{a} \mod$.
            \DontPrintSemicolon\;

        \nl Both $\A$ and $\B$ can compute the same private key $\kappa_{\A\B} \equiv \beta^{a} \mod$, 
            since $\beta^{a} = \alpha^{b}$.
            \DontPrintSemicolon\;

        \nl Assuming $\E$ intercepts either of $\alpha$ and $\beta$,
            it won't be able to retrive $\kappa_{\A\B}$ since $a$ and $b$ are unknown.
    \end{procedure}
\end{document}
